%% kapa: exam=1 round=36 year=2025 period=1 form=A subject=economics subject_ko=경제학원론 seq=2 number=42
%% source: Q-Net 감정평가사 기출문제 (cst00309) artlSeq=5231525 file=raw/1008/5231525/2025년 제36회 감정평가사 1차시험 1교시 문제지 원본.pdf
%% conversion: kordoc 4.24.51 -> tools/gen1.py (auto); review=proofread-sonnet; flags=none
\begin{question}{42}
소비자가 X재의 판매자일 때, 실물부존(real endowment) 모형에 관한 설명으로 옳지 않은 것은? (단, 주어진 소득으로 X재와 Y재만 양(+)의 소비를 한다.)
\begin{choices}
  \item 최적소비점은 예산선과 무차별곡선이 접하는 점이며 실물부존점과 일치하지 않을 수 있다.
  \item X재의 가격이 하락하고 효용이 증가하였다면, 소비자는 X재 구매자로 전환된다.
  \item Y재의 가격이 하락하면, 소비자의 효용은 증가한다.
  \item X재의 가격이 상승하고 Y재가 정상재라면, Y재의 소비는 증가한다.
  \item X재의 가격이 상승하고 X재가 열등재라면, X재의 소비는 증가한다.
\end{choices}
\end{question}
\answer{5}
